\section{Many-body response, screening, and conservation laws}
\label{sec:manybody}

\subsection{Why screening, polarization, and the vertex are not independent}

In Hedin's closed set of equations, the irreducible polarization $P$, dielectric
function $\epsilon$, screened interaction $W$, and reducible density response
$\chi$ satisfy schematically \citep{Hedin1965,Giustino2017}
\begin{align}
  P &= -\ii GG\Gamma_\rho, \\
  \epsilon &= 1-vP, \\
  W &= \epsilon^{-1}v, \\
  \chi &= P(1-vP)^{-1}.
  \label{eq:hedinresponse}
\end{align}
Here $\Gamma_\rho=-\delta G^{-1}/\delta V_{\mathrm{tot}}$ is a proper density
vertex in the stated convention.  Matrix products and integrations over internal
variables are implicit.  These relations answer an important conceptual question:
screening, the quasiparticle vertex, and static density response are not three
independent physical factors available for arbitrary multiplication.  The vertex
enters $P$; $P$ determines $\epsilon$; and $\chi$ follows from both.  Mixing proper
and screened vertex conventions can therefore count screening twice.

In a scalar electron gas, a convenient schematic form for the quasiparticle
coupling to an external ionic perturbation is
\begin{equation}
  g^{\mathrm{phys}}(\vq,\omega)
  = g^0(\vq)\,
    \frac{z\,\Gamma_\rho(\vq,\omega)}
         {\epsilon(\vq,\omega)},
  \label{eq:gphysical}
\end{equation}
where $z$ is the quasiparticle residue.  This equation is useful only after the
vertex convention and screening factor have been stated.  In this scalar
expression $z\Gamma_\rho$ includes the quasiparticle external-leg residue
appropriate to that convention.  It must not be identified without conversion
with the fixed-basis inverse-Green-function vertex of \cref{eq:dmb}.

\subsection{The Ward identity and its two limits}

Charge conservation and gauge invariance imply the Ward--Takahashi identity
\citep{Nambu1960}
\begin{equation}
  \ii\nu\,\Gamma^0(k+q,k)
  -\vq\cdot\bm{\Gamma}^{J}(k+q,k)
  = G^{-1}(k+q)-G^{-1}(k),
  \label{eq:ward}
\end{equation}
up to sign conventions for the Fourier transform and electric charge.
$\Gamma^0$ is the scalar charge vertex and $\bm{\Gamma}^{J}$ is the current
vertex.  The identity constrains the full four-vector vertex; the conclusion
depends on how $(\vq,\nu)$ approaches zero.

If one sets $\vq=0$ first and then takes $\nu\to0$, the dynamic uniform limit is
\begin{equation}
  \Gamma^0_{\omega}
  = \frac{\partial G^{-1}}{\partial(\ii\omega)}
  = 1-\frac{\partial\Sigma}{\partial(\ii\omega)}
  = z^{-1},
  \qquad z\Gamma^0_{\omega}=1.
  \label{eq:dynamicward}
\end{equation}
This is an exact cancellation for the conserved total charge in that limit.

If instead one sets $\nu=0$ first and then lets $\vq\to0$, the Ward identity
directly constrains the longitudinal current vertex.  The static scalar response
also contains the response of the self-energy to chemical potential and the
Landau interaction.  For an isotropic Fermi liquid, for example,
\begin{equation}
  \frac{\partial n}{\partial\mu}
  = \frac{N^*(0)}{1+F_0^s},
  \label{eq:compressibility}
\end{equation}
not simply $zN_0(0)$.  Therefore $z\Gamma_\rho=1$ is not a universal static
density-vertex identity.

For an acoustic phonon, $\omega_{\nu\vq}\simeq c_s|\vq|$ and generally
$v_F/c_s\gg1$, so
\begin{equation}
  \frac{v_F|\vq|}{\omega_{\nu\vq}}
  \simeq \frac{v_F}{c_s}\gg1.
  \label{eq:acousticstatic}
\end{equation}
This is the adiabatic or static side of the response.  The dynamic result in
\cref{eq:dynamicward} cannot by itself explain the success of ordinary DFPT for
acoustic phonons.

\subsection{The exact common-potential anchor}

There is nevertheless a clear physical picture in the strict uniform case.  If a
lattice displacement produces a potential $H'=Du\hat N$ that couples only to the
conserved total particle number, every many-electron state with $N$ electrons
shifts by $NDu$.  The energy required to add one electron therefore shifts by
exactly $Du$, irrespective of the strength of interactions inside the system.

Equivalently, a common scalar shift can be absorbed into the frequency origin,
\begin{equation}
  \Sigma_u(\omega)=\Sigma_0(\omega-Du),
  \qquad
  \partial_u\Sigma=-D\,\partial_\omega\Sigma.
  \label{eq:selfenergyshift}
\end{equation}
The shift of the quasiparticle pole is then
\begin{equation}
  z\bigl[D+\partial_u\Sigma\bigr]=D.
  \label{eq:protectedshift}
\end{equation}
This derivation explains the cancellation physically: an interaction cannot
renormalize a change of the energy zero.  It is an exact anchor at strict $\vq=0$
for a common charge potential, not a proof at finite momentum and not a statement
about orbital splitting, bond modulation, or other internal operators.

\begin{keypoint}[What conservation does for this project]
Conservation laws identify one protected full-system direction: a strict uniform
$\vq=0$ scalar potential coupled to total charge.  A projected
\texttt{global\_charge} component may use this anchor only after the original
unprojected perturbation is verified to be that full-space common shift.
Relative site/block identities (\texttt{site\_charge}), block-internal
traceless operators (\texttt{internal}), and inter-block operators
(\texttt{nonlocal}) are not protected by this argument.
\end{keypoint}
